% Unit 039 terjemahan Bahasa Indonesia dari chapter3.tex baris 1293--1586.
% Sumber: Methods of Algebra, Volume 2, commit
% 9a5803ff77dd3257484cb177f851a73770a59dd3 (CC BY 4.0).
% stable-unit-id: o014.aljabr2.chapter3.exercises-hochschild-homology-and-cohomology
% Cakupan: Bagian 3.8 lengkap.

% segment-id: o014.li2.chapter3.exercises-hochschild-homology-and-cohomology.h001
\section{Latihan: Homologi dan Kohomologi Hochschild}\label{sec:HH}
\hypertarget{o014.aljabr2.chapter3.exercises-hochschild-homology-and-cohomology}{ }

% segment-id: o014.li2.chapter3.exercises-hochschild-homology-and-cohomology.p001
Bagian ini kembali pada matematika konkret. Ambil $\Bbbk$ sebagai gelanggang
komutatif. Hasil kali tensor $\dotimes{\Bbbk}$ antara modul-$\Bbbk$ akan
disingkat menjadi $\otimes$, pangkat tensor ke-$n$ dari modul-$\Bbbk$ $M$
akan ditulis $M^{\otimes n}$, dan kita menetapkan
$M^{\otimes 0}:=\Bbbk$. Untuk aljabar-$\Bbbk$ $R$, seperti biasa, perkalian
kiri dan kanan oleh $\Bbbk$ pada setiap bimodul $(R,R)$ diasumsikan sama.
Karena itu, bimodul $(R,R)$ juga diidentifikasi dengan modul kiri
$R\otimes R^{\opp}$.

% segment-id: o014.li2.chapter3.exercises-hochschild-homology-and-cohomology.de001
\begin{definisi}[Kompleks bar]\label{def:bar-Hochschild}
	\index{kompleks bar (bar complex)}
	\index[sym1]{BR@$\mathsf{B}R$}
	Tetapkan suatu aljabar-$\Bbbk$ $R$. Untuk setiap
	$n\in\Z_{\geq 0}$, definisikan bimodul $(R,R)$
% segment-id: o014.li2.chapter3.exercises-hochschild-homology-and-cohomology.q001
	\begin{gather*}
		\mathsf{B}_n R := R \otimes R^{\otimes n} \otimes R = R^{\otimes (n+2)}, \\
		r(r_0 \otimes \cdots \otimes r_{n+1})r'
		= rr_0 \otimes \cdots \otimes r_{n+1} r',
	\end{gather*}
	dengan $r,r',r_0,\ldots,r_{n+1}\in R$. Untuk setiap $n\geq1$,
	definisikan homomorfisme bimodul
% segment-id: o014.li2.chapter3.exercises-hochschild-homology-and-cohomology.q002
	\begin{equation*}\begin{split}
		b_n: \mathsf{B}_n R & \to \mathsf{B}_{n-1} R, \\
		b_n \left( r_0 \otimes \cdots \otimes r_{n+1} \right)
		& := \sum_{k=0}^n (-1)^k
		\cdots \otimes r_k r_{k+1} \otimes \cdots .
	\end{split}\end{equation*}
	Kita menyebut
	$\mathsf{B}R:=\left(\mathsf{B}_nR,b_n\right)_{n\geq0}$ sebagai
	\emph{kompleks bar} dari $R$; ini adalah kompleks rantai dalam
	pengertian Catatan~\ref{rem:cochain-vs-chain}.
\end{definisi}

% segment-id: o014.li2.chapter3.exercises-hochschild-homology-and-cohomology.p002
Notasi klasik untuk
$r_0\otimes\cdots\otimes r_{n+1}\in\mathsf{B}_nR$ ialah
$(r_0|\cdots|r_{n+1})$; dari sinilah nama ``bar'' berasal. Homomorfisme $b$
bekerja dengan menghapus satu garis pemisah dengan segala cara yang mungkin,
lalu menjumlahkan hasilnya dengan tanda berselang-seling. Dari sini mudah
dilihat bahwa $b^2=0$. Memang,
$b^2(r_0|\cdots|r_{n+1})$ merupakan kombinasi linear dari unsur berbentuk
% segment-id: o014.li2.chapter3.exercises-hochschild-homology-and-cohomology.q003
\begin{gather*}
	(\cdots |r_{h-1} r_h| \cdots |r_{k-1} r_k| \cdots )
	\quad \text{atau} \quad
	( \cdots | r_{h-1} r_h r_{h+1} | \cdots ).
\end{gather*}
Untuk memperoleh suku semacam itu dari
$(\cdots|r_{h-1}|\cdots|r_k|\cdots)$ dengan menghapus garis pemisah,
terdapat tepat dua cara, dan tanda keduanya saling meniadakan. Jadi,
$\mathsf{B}R$ memang merupakan kompleks rantai.

% segment-id: o014.li2.chapter3.exercises-hochschild-homology-and-cohomology.p003
Sekarang augmentasikan kompleks bar seperti pada diagram berikut dan tuliskan
hasilnya sebagai $\mathsf{B}'R$. Morfisme $b_0$ pada diagram juga disebut
homomorfisme augmentasi.
% segment-id: o014.li2.chapter3.exercises-hochschild-homology-and-cohomology.g001
\[\begin{tikzcd}[row sep=small]
	\cdots \arrow[r, "b_2"] & \mathsf{B}_1 R \arrow[r, "b_1"]
	& \mathsf{B}_0 R \arrow[r, "b_0"] & (\mathsf{B}_{-1} R := R) \arrow[r] & 0 \\
	& & (r_0 | r_1) \arrow[mapsto, r]
	\arrow[phantom, u, "\in" description, sloped]
	& r_0 r_1 \arrow[phantom, u, "\in" description, sloped] &
\end{tikzcd}\]\index[sym1]{B'R@$\mathsf{B}'R$}

% segment-id: o014.li2.chapter3.exercises-hochschild-homology-and-cohomology.le001
\begin{lema}\label{prop:HH-bar-exactness}
	Kompleks rantai teraugmentasi $\mathsf{B}'R$ bersifat eksak. Lebih tepatnya,
	jika $\mathsf{B}'R$ dipandang sebagai kompleks rantai modul-$\Bbbk$, maka
	$\identity_{\mathsf{B}'R}$ homotopik nol.
\end{lema}

% segment-id: o014.li2.chapter3.exercises-hochschild-homology-and-cohomology.pf001
\begin{bukti}
	Kita akan membangun suatu keluarga homomorfisme modul-$\Bbbk$
	$h_n:\mathsf{B}'_nR\to\mathsf{B}'_{n+1}R$, untuk $n\geq-1$, sedemikian
	sehingga
	$b_{n+1}h_n+h_{n-1}b_n=\identity_{\mathsf{B}'_nR}$ (tentu dengan konvensi
	$h_{-2}=0$ dan $b_{-1}=0$). Ini akan membuktikan bahwa
	$\identity_{\mathsf{B}'R}$ homotopik nol. Secara konkret, ambil
% segment-id: o014.li2.chapter3.exercises-hochschild-homology-and-cohomology.q004
	\[ h_n(r_0 | \cdots | r_{n+1})
		= (1_R | r_0 | \cdots | r_{n+1}), \]
	lalu perluas secara linear pada $\mathsf{B}'_nR$. Untuk setiap $n\geq0$
	dan $r_0,\ldots,r_{n+1}\in R$, kita mempunyai
% segment-id: o014.li2.chapter3.exercises-hochschild-homology-and-cohomology.q005
	\begin{align*}
		b_{n+1} h_n \left(r_0 | \cdots | r_{n+1} \right)
		& = (r_0 | \cdots | r_{n+1})
		+ \sum_{k=0}^n (-1)^{k+1}
		(1_R | \cdots | r_k r_{k+1} | \cdots), \\
		h_{n-1} b_n \left(r_0 | \cdots | r_{n+1} \right)
		& = \sum_{k=0}^n (-1)^k
		(1_R | \cdots | r_k r_{k+1} | \cdots).
	\end{align*}
	Dengan demikian, segera diperoleh
% segment-id: o014.li2.chapter3.exercises-hochschild-homology-and-cohomology.q006
	\[ (b_{n+1} h_n + h_{n-1} b_n)(r_0 | \cdots | r_{n+1})
		= (r_0 | \cdots | r_{n+1}). \]
	Karena $b_0h_{-1}(r)=b_0(1_R|r)=r$, persamaan ini juga berlaku secara
	langsung untuk $n=-1$.
\end{bukti}

% segment-id: o014.li2.chapter3.exercises-hochschild-homology-and-cohomology.p004
Sekilas, definisi kompleks $\mathsf{B}R$ dan $\mathsf{B}'R$ tampak seperti
rekaan yang tidak terduga.
Contoh~\sourcecrossref{eg:HH-comonad}{pada pembahasan Bab~8} nanti akan
memandang kompleks-kompleks rantai ini dari sudut pandang komonad dan
memberikan penjelasan yang alami.

% segment-id: o014.li2.chapter3.exercises-hochschild-homology-and-cohomology.cn001
\begin{konvensi}\label{con:Re-R-R}
	\index[sym1]{R-e@$R^e$}
	Tetapkan $R^e:=R\otimes R^{\opp}$. Untuk setiap bimodul $(R,R)$ $M$,
	termasuk kasus khusus $M=R$, mulai sekarang:
% segment-id: o014.li2.chapter3.exercises-hochschild-homology-and-cohomology.l001
	\begin{compactitem}
% segment-id: o014.li2.chapter3.exercises-hochschild-homology-and-cohomology.i001
		\item pandang $M$ sebagai modul kiri $R^e$ melalui
		$(r\otimes r')m=rmr'$;
% segment-id: o014.li2.chapter3.exercises-hochschild-homology-and-cohomology.i002
		\item pandang $M$ sebagai modul kanan $R^e$ melalui
		$m(r\otimes r'):=r'mr$.
	\end{compactitem}
\end{konvensi}

% segment-id: o014.li2.chapter3.exercises-hochschild-homology-and-cohomology.p005
Dengan konvensi ini, untuk setiap $M$ kita dapat mendefinisikan kompleks
rantai modul-$\Bbbk$ $M\dotimes{R^e}\mathsf{B}R$ dan kompleks
$\Hom_{R^e}\left(\mathsf{B}R,M\right)$.

% segment-id: o014.li2.chapter3.exercises-hochschild-homology-and-cohomology.de002
\begin{definisi}[G.\ Hochschild]\label{def:HH}
	\index{homologi dan kohomologi Hochschild}
	\index[sym1]{HH@$\HHm_n$, $\HHm^n$}
	Untuk setiap $n$, definisikan funktor-funktor $\Bbbk$-linear
	$\HHm_n,\HHm^n:(R,R)\dcate{Mod}\rightrightarrows\Bbbk\dcate{Mod}$
	sebagai berikut:
% segment-id: o014.li2.chapter3.exercises-hochschild-homology-and-cohomology.q007
	\[\begin{array}{cl}
		\text{Homologi Hochschild}
		& \HHm_n(M) := \Hm_n\left(M\dotimes{R^e}\mathsf{B}R\right), \\
		\text{Kohomologi Hochschild}
		& \HHm^n(M) := \Hm^n\left(\Hom_{R^e}(\mathsf{B}R,M)\right).
	\end{array}\]
	Dalam kasus khusus $M=R$, objek-objek ini masing-masing disebut
	\emph{homologi Hochschild} $\HHm_n(R)$ dan \emph{kohomologi Hochschild}
	$\HHm^n(R)$ dari $R$.
\end{definisi}

% segment-id: o014.li2.chapter3.exercises-hochschild-homology-and-cohomology.p006
Untuk menyederhanakan deskripsi $\HHm_n(M)$ dan $\HHm^n(M)$, perkenalkan dua
kompleks Hochschild berikut:
% segment-id: o014.li2.chapter3.exercises-hochschild-homology-and-cohomology.q008
\begin{equation}\label{eqn:HH-cplx}\begin{aligned}
	\index[sym1]{CRM@$C_\bullet(R, M)$, $C^\bullet(R, M)$}
	C_\bullet(R,M) & := \left[
	\cdots \to M \otimes R^{\otimes n} \xrightarrow{d_n} \cdots
	\xrightarrow{d_2} M \otimes R \xrightarrow{d_1} M \right], \\
	C^\bullet(R,M) & := \left[
	M \xrightarrow{d^0} \Hom_{\Bbbk}(R, M) \xrightarrow{d^1} \cdots
	\to \Hom_{\Bbbk}(R^{\otimes n}, M) \xrightarrow{d^n} \cdots \right].
\end{aligned}\end{equation}
\index{kompleks Hochschild (Hochschild complex)}
Derajatnya masing-masing ialah $\ldots,2,1,0$ dan $0,1,2,\ldots$; semua suku
lainnya diisi dengan $0$. Homomorfisme $d_n$ dan $d^n$ didefinisikan sebagai
berikut.
% segment-id: o014.li2.chapter3.exercises-hochschild-homology-and-cohomology.l002
\begin{enumerate}
% segment-id: o014.li2.chapter3.exercises-hochschild-homology-and-cohomology.i003
	\item Dengan notasi bar, tetap tuliskan unsur
	$m\otimes r_1\otimes\cdots\otimes r_n$ dari $M\otimes R^{\otimes n}$
	sebagai $(m|r_1|\cdots|r_n)$. Tetapkan
% segment-id: o014.li2.chapter3.exercises-hochschild-homology-and-cohomology.q009
	\begin{multline}\label{eqn:HH-cplx-d}
		d_n(m|r_1|\cdots|r_n)=\\
		\underbracket{(mr_1|r_2|\cdots|r_n)
		+\sum_{k=1}^{n-1}(-1)^k
		(m|\cdots|r_kr_{k+1}|\cdots)}_{=:d'_n(m|r_1|\cdots|r_n)}
		+(-1)^n(r_nm|r_1|\cdots|r_{n-1}).
	\end{multline}
% segment-id: o014.li2.chapter3.exercises-hochschild-homology-and-cohomology.i004
	\item Identifikasikan unsur $\Hom_{\Bbbk}(R^{\otimes n},M)$ dengan
	pemetaan $\Bbbk$-multilinear $R^n\to M$ (lihat
	\cite[\S 7.5]{Li1}). Tetapkan
% segment-id: o014.li2.chapter3.exercises-hochschild-homology-and-cohomology.q010
	\begin{multline*}
		(d^nf)(r_1,\ldots,r_{n+1})=\\
		r_1f(r_2,\ldots,r_{n+1})
		+\sum_{k=1}^n(-1)^kf(\ldots,r_kr_{k+1},\ldots)
		+(-1)^{n+1}f(r_1,\ldots,r_n)r_{n+1}.
	\end{multline*}
\end{enumerate}

% segment-id: o014.li2.chapter3.exercises-hochschild-homology-and-cohomology.p007
Untuk setiap $n\in\Z_{\geq0}$, kita mempunyai isomorfisme
% segment-id: o014.li2.chapter3.exercises-hochschild-homology-and-cohomology.g002
\[\begin{tikzcd}[row sep=tiny]
	M \dotimes{R^e} \mathsf{B}_n R \arrow[leftrightarrow, r, "\sim"]
	& M \otimes R^{\otimes n} \\
	m \otimes (r_0 \otimes \cdots \otimes r_{n+1}) \arrow[mapsto, r]
	& (r_{n+1} m r_0 | r_1 | \cdots | r_n) \\
	m \otimes (1_R \otimes r_1 \otimes \cdots \otimes r_n \otimes 1_R )
	& (m | r_1 | \cdots | r_n) \arrow[mapsto, l]
\end{tikzcd}\]
\footnote{Catatan penerjemah (O014-C044): sumber menambahkan komponen akhir
$|1_R$ pada
citra di baris tengah, meskipun kodomainnya $M\otimes R^{\otimes n}$; komponen
berlebih itu dihapus di sini.}
dan
% segment-id: o014.li2.chapter3.exercises-hochschild-homology-and-cohomology.g003
\[\begin{tikzcd}[row sep=tiny, column sep=small]
	\Hom_{R^e}(\mathsf{B}_nR, M) \arrow[leftrightarrow, r, "\sim"]
	& \Hom_{\Bbbk}(R^{\otimes n}, M) \\
	\varphi \arrow[mapsto, r]
	& {\left[(r_1,\ldots,r_n)\mapsto
		\varphi(1_R,r_1,\ldots,r_n,1_R)\right]} \\
	{\left[r_0\otimes\cdots\otimes r_{n+1}\mapsto
		r_0f(r_1,\ldots,r_n)r_{n+1}\right]}
	& f \arrow[mapsto, l] .
\end{tikzcd}\]
Perhitungan singkat menunjukkan bahwa, melalui isomorfisme ini,
$\identity_M\otimes b_n$ bersesuaian dengan $d_n$, sedangkan $b_{n+1}^*$
bersesuaian dengan $d^n$.\footnote{Catatan penerjemah (O014-C045): sumber menulis
$b_n^*$, yang berawal pada $\Hom_{R^e}(\mathsf{B}_{n-1}R,M)$; diferensial
berderajat $n$ berasal dari $b_{n+1}:\mathsf{B}_{n+1}R\to\mathsf{B}_nR$.}
Jadi, isomorfisme-isomorfisme di atas merupakan isomorfisme kompleks dan
memberikan hasil berikut. Untuk menampilkan kembali parameter $R$ secara
eksplisit, tuliskan
$\HHm_n(R,M):=\HHm_n(M)$ dan $\HHm^n(R,M):=\HHm^n(M)$.
\footnote{Catatan penerjemah (O014-C046): pada rumus berikut sumber beralih ke notasi
dua argumen tanpa mendefinisikannya; kedua singkatan ini menjadikan peralihan
tersebut eksplisit.} Dengan notasi ini,
% segment-id: o014.li2.chapter3.exercises-hochschild-homology-and-cohomology.q011
\[ \HHm^n(R, M) \simeq \Hm^n(C^\bullet(R, M)), \quad
	\HHm_n(R, M) \simeq \Hm_n(C_\bullet(R, M)). \]
Pembaca yang baru mempelajari topik ini dianjurkan untuk memeriksa
perhitungannya secara terperinci.

% segment-id: o014.li2.chapter3.exercises-hochschild-homology-and-cohomology.re001
\begin{catatan}
	Misalkan $R$ komutatif. Setiap modul-$R$ $M$ menjadi bimodul $(R,R)$
	melalui $rmr':=rr'm$. Dalam hal ini, kompleks $C_\bullet(R,M)$ (atau
	$C^\bullet(R,M)$)\footnote{Catatan penerjemah (O014-C047): sumber menukar kedua
	argumen menjadi $C_\bullet(M,R)$ dan $C^\bullet(M,R)$ pada kalimat ini;
	notasi di sini diselaraskan dengan definisi sebelumnya.}
	merupakan kompleks modul-$R$ jika kita menetapkan
% segment-id: o014.li2.chapter3.exercises-hochschild-homology-and-cohomology.q012
	\[ r\cdot(m|r_1|\cdots|r_n)=(rm|r_1|\cdots|r_n)
		\quad\text{atau}\quad
		(r\cdot f)(r_1,\ldots,r_n)=rf(r_1,\ldots,r_n). \]
	Karena itu, $\HHm_n(M)$ dan $\HHm^n(M)$ keduanya memperoleh struktur
	modul-$R$. Hal ini tentu juga dapat dibuktikan pada tingkat kompleks bar.
\end{catatan}

% segment-id: o014.li2.chapter3.exercises-hochschild-homology-and-cohomology.ex001
\begin{contoh}
	Ambil $M=R=\Bbbk$. Dari konstruksi di atas langsung terlihat bahwa
	$d_1,d_2,\ldots$ berturut-turut adalah
	$0,\identity,0,\identity,\ldots$. Jadi
	$\HHm_0(\Bbbk)=\Bbbk$ dan $\HHm_{\geq1}(\Bbbk)=0$. Argumen serupa
	memberikan $\HHm^0(\Bbbk)=\Bbbk$ dan $\HHm^{\geq1}(\Bbbk)=0$.
\end{contoh}

% segment-id: o014.li2.chapter3.exercises-hochschild-homology-and-cohomology.p008
Untuk derajat $n$ yang lebih tinggi, menghitung $\HHm^n(M)$ dan $\HHm_n(M)$
dari definisi awal pada umumnya sulit.
Contoh~\sourcecrossref{eg:HH-Tor}{pada \S~3.14} dan
\sourcecrossref{eg:HH-Ext}{pada \S~3.14} nanti masing-masing akan menunjukkan
bahwa
% segment-id: o014.li2.chapter3.exercises-hochschild-homology-and-cohomology.q013
\[ \HHm_n(M) \simeq \Tor^{R^e}_n(M, R), \quad
	\HHm^n(M) \simeq \Ext_{R^e}^n(R, M), \]
dengan syarat bahwa $R$, sebagai modul-$\Bbbk$, masing-masing datar dan
proyektif. Pada saat itu, beberapa kasus khusus $\HHm_n$ dan $\HHm^n$ dapat
dihitung dengan kompleks yang lebih sederhana. Untuk $R$ umum, pada bagian
latihan \sourcecrossref{sec:simplicial}{Bab~8} kita akan menjelaskan
$\HHm_n$ dan $\HHm^n$
sebagai $\Tor$ relatif dan $\Ext$ relatif.

% segment-id: o014.li2.chapter3.exercises-hochschild-homology-and-cohomology.ex002
\begin{contoh}[Derajat nol: pusat dan kopusat]
	\index{pusat}
	\index{kopusat (cocenter)}
	Pertama, tinjau $\HHm^0(M)$. Menurut definisi,
	$d^0:M=C^0(R,M)\to C^1(R,M)=\Hom_{\Bbbk}(R,M)$ memetakan $m$ ke
	$[r\mapsto rm-mr]$. Karena itu,
% segment-id: o014.li2.chapter3.exercises-hochschild-homology-and-cohomology.q014
	\[ \HHm^0(M)=\left\{m\in M:\forall r\in R,\;rm=mr\right\}. \]
	Ruas kanan secara wajar dapat disebut \emph{pusat} dari $M$; ketika
	$M=R$, ini tidak lain daripada pusat dalam pengertian teori gelanggang.

	Selanjutnya, tinjau $\HHm_0(M)$. Tuliskan $[M,R]$ untuk submodul-$\Bbbk$
	dari $M$ yang dibangun oleh semua unsur berbentuk $mr-rm$, dengan
	$m\in M$ dan $r\in R$. Karena $d_1(m|r)=mr-rm$, kita memperoleh
% segment-id: o014.li2.chapter3.exercises-hochschild-homology-and-cohomology.q015
	\[ \Image\left[d_1:M\otimes R\to M\right]=[M,R],
		\quad \HHm_0(M)=M/[M,R]. \]
	Khususnya, kasus $M=R$ memberikan submodul $[R,R]$ dari $R$ yang
	dibangun oleh semua $[r',r]:=r'r-rr'$. Modul-$\Bbbk$ yang bersesuaian,
	$R/[R,R]$, disebut \emph{kopusat} dari $R$. Setiap homomorfisme
	modul-$\Bbbk$ $\varphi:R\to N$ yang memenuhi
	$\varphi(rr')=\varphi(r'r)$, yakni yang mempunyai sifat menyerupai
	``jejak'', berfaktor secara unik melalui
	$R/[R,R]=\HHm_0(R)$.
\end{contoh}

% segment-id: o014.li2.chapter3.exercises-hochschild-homology-and-cohomology.ex003
\begin{contoh}[Derajat satu: derivasi]\label{eg:HH1}
	\index[sym1]{Der@$\mathrm{Der}_{\Bbbk}(R, M)$}
	\index[sym1]{InnDer@$\mathrm{Inn}_{\Bbbk}(R, M)$}
	Sekarang tinjau $\HHm^1(M)$. Tuliskan $[r,m]:=rm-mr$. Maka
% segment-id: o014.li2.chapter3.exercises-hochschild-homology-and-cohomology.q016
	\begin{align*}
		\Ker(d^1)
		&=\left\{D\in\Hom_{\Bbbk}(R,M):
		\forall r_1,r_2\in R,\;
		r_1D(r_2)-D(r_1r_2)+D(r_1)r_2=0\right\},\\
		\Image(d^0)
		&=\left\{[\mathord{\cdot},m]\in\Hom_{\Bbbk}(R,M):m\in M\right\}.
	\end{align*}
	Syarat yang mendefinisikan $\Ker(d^1)$ dapat ditulis ulang sebagai
	aturan Leibniz
	$D(r_1r_2)=r_1D(r_2)+D(r_1)r_2$. Homomorfisme modul-$\Bbbk$ $D$ dengan
	sifat ini dipandang sebagai operasi derivasi bernilai di $M$. Modul
	$\Bbbk$ yang dibentuknya ditulis $\Der_{\Bbbk}(R,M)$, sedangkan
	unsur-unsur berbentuk $[\mathord{\cdot},m]$ membentuk submodul
	$\mathrm{Inn}_{\Bbbk}(R,M)$. Jadi, modul hasil bagi
% segment-id: o014.li2.chapter3.exercises-hochschild-homology-and-cohomology.q017
	\[ \HHm^1(M)=\Der_{\Bbbk}(R,M)/\mathrm{Inn}_{\Bbbk}(R,M) \]
	mengklasifikasikan semua operasi derivasi pada $R$ yang bernilai di $M$,
	modulo $\mathrm{Inn}_{\Bbbk}(R,M)$.

	Sekarang anggap $R$ komutatif untuk menafsirkan $\HHm_1(M)$. Kita
	memerlukan sedikit persiapan. Definisikan modul bebas-$R$
	$\bigoplus_{r\in R}R\widetilde{\dd r}$ dengan basis simbol-simbol
	$\widetilde{\dd r}$, lalu definisikan submodul $N$ yang dibangun oleh
	unsur-unsur berikut:
% segment-id: o014.li2.chapter3.exercises-hochschild-homology-and-cohomology.q018
	\[ \widetilde{\dd(r+r')}-\widetilde{\dd r}-\widetilde{\dd r'},
		\quad \widetilde{\dd tr}-t\widetilde{\dd r},
		\quad \widetilde{\dd rr'}-r\widetilde{\dd r'}-r'\widetilde{\dd r}, \]
	dengan $r,r'\in R$ dan $t\in\Bbbk$. Dengan demikian, definisikan
	modul-$R$
% segment-id: o014.li2.chapter3.exercises-hochschild-homology-and-cohomology.q019
	\begin{align*}
		\Omega_{R|\Bbbk}
		&:=\bigoplus_{r\in R}R\widetilde{\dd r}\bigg/N\\
		&=\sum_{r\in R}R\dd r,\qquad
		\dd r:=\text{citra dari }\widetilde{\dd r}.
	\end{align*}
	\index{diferensial Kähler (Kähler differentials)}
	\index[sym1]{OmegaRk@$\Omega_{R\vert\Bbbk}$}
	Modul ini disebut \emph{modul diferensial Kähler} dari aljabar-$\Bbbk$
	$R$ dan dicirikan oleh sifat universal berikut:
% segment-id: o014.li2.chapter3.exercises-hochschild-homology-and-cohomology.g004
	\[\begin{tikzcd}[row sep=tiny]
		\Hom_R(\Omega_{R|\Bbbk}, M) \arrow[leftrightarrow, r, "\sim"]
		& \Der_{\Bbbk}(R, M) \\
		\varphi \arrow[mapsto, r] & {\left[r\mapsto\varphi(\dd r)\right]} \\
		{\left[\dd r\mapsto D(r)\right]} & D \arrow[mapsto, l].
	\end{tikzcd}
	\qquad M:\text{modul-$R$}. \]
	Pembaca diminta memeriksanya secara langsung. Isomorfisme ini menunjukkan
	bahwa $R\to\Omega_{R|\Bbbk}$ yang diberikan oleh $r\mapsto\dd r$
	(bersesuaian dengan $\varphi=\identity$) merupakan ``derivasi universal''.
	Topik yang berkaitan dengan hal ini semestinya menjadi bagian dari teori
	gelanggang komutatif; kita cukup mencatatnya di sini.

	Kembali ke $\HHm_1(M)$, tetap anggap $R$ komutatif. Misalkan $M$ adalah
	modul-$R$ yang dijadikan bimodul $(R,R)$ melalui $rmr':=(rr')m$.
	Langsung terlihat bahwa $d_1:M\otimes R\to M$ adalah $0$, sedangkan citra
	$d_2:M\otimes R^{\otimes2}\to M\otimes R$ dibangun oleh unsur-unsur
	berbentuk $(rm|r')-(m|rr')+(r'm|r)$. Maka kita memperoleh homomorfisme
	modul-$\Bbbk$ dua arah
% segment-id: o014.li2.chapter3.exercises-hochschild-homology-and-cohomology.g005
	\[\begin{tikzcd}[row sep=tiny]
		(M \otimes R) / \Image(d_2) \arrow[shift right, r]
		& M \dotimes{R} \Omega_{R|\Bbbk} \arrow[shift right, l] \\
		(m|r) + \Image(d_2) \arrow[mapsto, r] & m \otimes \dd r \\
		(r'm|r) + \Image(d_2)
		& m \otimes r' \dd r \arrow[mapsto, l] .
	\end{tikzcd}\]
	Berdasarkan deskripsi $\Image(d_2)$ dan $\Omega_{R|\Bbbk}$, pemeriksaan
	rutin menunjukkan bahwa homomorfisme-homomorfisme tersebut terdefinisi
	dengan baik dan saling invers. Bahkan, keduanya merupakan isomorfisme
	modul-$R$. Jadi, ketika $R$ komutatif, kita memperoleh
	$\HHm_1(M)\simeq M\dotimes{R}\Omega_{R|\Bbbk}$. Khususnya,
	$\HHm_1(R)\simeq\Omega_{R|\Bbbk}$.
\end{contoh}

% segment-id: o014.li2.chapter3.exercises-hochschild-homology-and-cohomology.p009
Homologi dan kohomologi Hochschild memiliki kandungan yang kaya; hubungannya
dengan derivasi dan bentuk diferensial bukanlah suatu kebetulan. Bagian latihan
akan memberikan lebih banyak penafsiran terhadap kohomologi Hochschild.

% segment-id: o014.li2.chapter3.exercises-hochschild-homology-and-cohomology.p010
Kembali ke kompleks rantai Hochschild. Dengan memperhatikan aksi $R^e$ pada
$\mathsf{B}_nR$ dan pada $M$, gagasan intuitifnya ialah menyusun unsur
$(m|r_1|\cdots|r_n)$ dari $C_n(R,M)$ dalam bentuk melingkar:
% segment-id: o014.li2.chapter3.exercises-hochschild-homology-and-cohomology.f001
\begin{center}\begin{tikzpicture}[scale=1.2]
		\node[rotate=30] at (120:1) {$r_n$};
		\node[rotate = 0]  at (90:1) {$m$};
		\node[rotate=-15] at (60:1) {$r_1$};
		\node[rotate=-45]  at (30:1) {$r_2$};
		\draw (135:0.8) -- (135:1.2);
		\draw (105:0.8) -- (105:1.2);
		\draw (75:0.8) -- (75:1.2);
		\draw(45:0.8) -- (45:1.2);
		\draw[loosely dotted, thick]
		(10:1) arc [start angle=10, end angle=-210, radius=1];
\end{tikzpicture}\end{center}
Dengan demikian, $d_n(m|r_1|\cdots|r_n)$ diperoleh dengan menghapus garis
pemisah dengan $n+1$ cara dan menjumlahkan hasilnya dengan tanda
berselang-seling. Dalam kasus khusus $M=R$, gambar tersebut mempunyai simetri
rotasi yang jelas. Dari sudut pandang aljabar, untuk setiap
$n\in\Z_{\geq0}$ kita definisikan homomorfisme modul-$\Bbbk$
% segment-id: o014.li2.chapter3.exercises-hochschild-homology-and-cohomology.g006
\[\begin{tikzcd}[row sep=tiny]
	t_n: R^{\otimes (n+1)} \arrow[r] & R^{\otimes (n+1)} \\
	(r_0 | \cdots | r_n) \arrow[mapsto, r]
	& (-1)^n \cdot (r_n | r_0 | \cdots | r_{n-1}) ;
\end{tikzcd}\]
akibatnya, $t_n^{n+1}=\identity$. Homomorfisme ini mewujudkan simetri tersebut.

% segment-id: o014.li2.chapter3.exercises-hochschild-homology-and-cohomology.p011
Simetri rotasi mengarah pada teori \emph{homologi siklik}. Secara historis,
sekurang-kurangnya terdapat dua motivasi untuk mempelajari homologi siklik.
Pertama, mencari analog teori de Rham dalam situasi nonkomutatif. Kedua,
mempelajari dan menerapkan teori K, termasuk berbagai perumuman teorema
indeks. Bagian ini menggunakan homologi dan kohomologi Hochschild serta
homologi siklik hanya sebagai sarana ringan untuk membiasakan
pembaca dengan operasi pada kompleks; pembahasannya sekadar pengantar.
Pembaca yang ingin mendalaminya dapat melihat monograf seperti
\cite{Lo98, Wi19}.

% segment-id: o014.li2.chapter3.exercises-hochschild-homology-and-cohomology.p012
Semua kompleks (atau bikompleks) yang disebutkan selanjutnya secara bawaan
merupakan kompleks rantai (atau bikompleks rantai).

% segment-id: o014.li2.chapter3.exercises-hochschild-homology-and-cohomology.p013
Untuk setiap $n\geq0$, tetapkan
$N_n:=\identity+t_n+\cdots+(t_n)^n
\in\End_{\Bbbk}(R^{\otimes(n+1)})$. Definisikan \emph{bikompleks siklik}
$\mathrm{CC}(R)=(\mathrm{CC}(R)_{p,q})_{(p,q)\in\Z^2}$ sebagai bikompleks
berikut.
\index{bikompleks siklik (cyclic double complex)}
\index[sym1]{CCR@$\mathrm{CC}(R)$}
% segment-id: o014.li2.chapter3.exercises-hochschild-homology-and-cohomology.g007
\[\begin{tikzcd}[
	/tikz/execute at end picture={
		\node[rectangle, draw, fit=(NW) (SE),
		inner xsep=1.5em, inner ysep=1.5em] {};
	}]
	\vdots & |[alias=NW]| & \vdots \arrow[d, "{d_{q+1}}"]
	& \vdots \arrow[d, "{d'_{q+1}}"] & \\
	q & \cdots & R^{\otimes (q+1)} \arrow[d, "{d_q}"] \arrow[l, "{N_q}"]
	& R^{\otimes (q+1)} \arrow[d, "{d'_q}"]
	\arrow[l, "{\identity - t_q}"] & \cdots \arrow[l] \\
	q-1 & \cdots & R^{\otimes q} \arrow[d] \arrow[l, "{N_{q-1}}"]
	& R^{\otimes q} \arrow[d] \arrow[l, "{\identity - t_{q-1}}"]
	& \cdots \arrow[l] \\
	\vdots & & \vdots & \vdots & |[alias=SE]| \\
	& \cdots & p: \text{genap} & p+1: \text{ganjil} & \cdots
\end{tikzcd}\]
Semua sukunya merupakan modul-$\Bbbk$, dan semua suku dengan $q<0$
ditetapkan sebagai $0$. Morfisme $d_q$ dan $d'_q$ didefinisikan dalam
\eqref{eqn:HH-cplx-d}. Perhatikan bahwa:
% segment-id: o014.li2.chapter3.exercises-hochschild-homology-and-cohomology.l003
\begin{compactitem}
% segment-id: o014.li2.chapter3.exercises-hochschild-homology-and-cohomology.i005
	\item kolom genap adalah kompleks rantai Hochschild
	$C_\bullet(R,R)$ dari \eqref{eqn:HH-cplx};
% segment-id: o014.li2.chapter3.exercises-hochschild-homology-and-cohomology.i006
	\item mudah dibuktikan bahwa kolom ganjil adalah kompleks bar
	teraugmentasi $\mathsf{B}'R$, dengan indeks digeser sehingga dimulai
	pada derajat $0$. Karena itu, Lema~\ref{prop:HH-bar-exactness}
	menyiratkan bahwa semua kolom ganjil bersifat eksak; bahkan, kolom-kolom
	itu adalah $0$ dalam $\cate{K}(\Bbbk\dcate{Mod})$.
\end{compactitem}

% segment-id: o014.li2.chapter3.exercises-hochschild-homology-and-cohomology.p014
Untuk menunjukkan bahwa $\mathrm{CC}(R)$ memang merupakan bikompleks, kita
memerlukan pengamatan berikut. Argumennya elementer dan menarik, tanpa
kesulitan hakiki, sehingga diserahkan sebagai latihan dalam bab ini.

% segment-id: o014.li2.chapter3.exercises-hochschild-homology-and-cohomology.le002
\begin{lema}\label{prop:CC-cplx}
	Misalkan $R$ adalah aljabar-$\Bbbk$. Untuk setiap
	$q\in\Z_{\geq1}$, berlaku persamaan berikut dalam
	$\Hom_{\Bbbk}\left(R^{\otimes(q+1)},R^{\otimes q}\right)$:
% segment-id: o014.li2.chapter3.exercises-hochschild-homology-and-cohomology.q020
	\[ d_q(\identity-t_q)=(\identity-t_{q-1})d'_q,
		\quad d'_qN_q=N_{q-1}d_q. \]
\end{lema}

% segment-id: o014.li2.chapter3.exercises-hochschild-homology-and-cohomology.p015
Untuk setiap bikompleks $C=(C_{i,j})_{(i,j)\in\Z^2}$ dan $m\in\Z$,
definisikan pergeseran horizontalnya\footnote{Konvensi buku ini memakai
subskrip $\mathrm{I}$ untuk operasi horizontal dan $\mathrm{II}$ untuk
operasi vertikal.} $C_{\mathrm{I}}[m]$ sebagai bikompleks
$(C_{m+i,j})_{i,j}$. Untuk setiap $p\in\Z$, definisikan funktor pemenggalan
kasar horizontal (\textit{brutal truncation})
$\sigma_{\mathrm{I},\leq p}C$ (atau
$\sigma_{\mathrm{I},\geq p}C$) dengan mengganti $C_{i,j}$ yang memenuhi
$i>p$ (atau $i<p$) dengan $0$ dan membiarkan semua suku lain tetap.
Dengan demikian, kita mempunyai barisan eksak pendek bikompleks
(perhatikan urutannya):
% segment-id: o014.li2.chapter3.exercises-hochschild-homology-and-cohomology.q021
\begin{equation}\label{eqn:homology-I-brute-truncation}
	0 \to \sigma_{\mathrm{I}, \leq p} C \to C
	\to \sigma_{\mathrm{I}, \geq p+1} C \to 0.
\end{equation}
Selain itu, untuk setiap $a\leq b$, definisikan
$\sigma_{\mathrm{I},[a,b]}:=
\sigma_{\mathrm{I},\leq b}\sigma_{\mathrm{I},\geq a}
=\sigma_{\mathrm{I},\geq a}\sigma_{\mathrm{I},\leq b}$.


% segment-id: o014.li2.chapter3.exercises-hochschild-homology-and-cohomology.p016
Terapkan funktor-funktor tersebut pada $\mathrm{CC}(R)$, lalu ambil kompleks
total $\tot_{\Pi}$ (Definisi~\ref{def:total-cplx}); kita sampai pada definisi
berikut.

% segment-id: o014.li2.chapter3.exercises-hochschild-homology-and-cohomology.de003
\begin{definisi}[B.\ Feigin, B.\ Tsygan; A.\ Connes]\label{def:HC}
	\index{homologi siklik (cyclic homology)}
	\index{homologi siklik periodik (periodic cyclic homology)}
	\index[sym1]{HPHC@$\mathrm{HP}_n$, $\mathrm{HC}_n$, $\mathrm{HP}^n$, $\mathrm{HC}^n$}
	Untuk aljabar-$\Bbbk$ $R$ dan setiap $n\in\Z$, definisikan
% segment-id: o014.li2.chapter3.exercises-hochschild-homology-and-cohomology.q022
	\begin{align*}
		\mathrm{HP}_n(R)
		&:=\Hm_n\left(\tot_{\Pi}(\mathrm{CC}(R))\right)
		&\text{(homologi siklik periodik)},\\
		\mathrm{HC}_n(R)
		&:=\Hm_n\left(
		\tot\left(\sigma_{\mathrm{I},\geq0}\mathrm{CC}(R)\right)\right)
		&\text{(homologi siklik)}.
%		\mathrm{HC}^-_n(R) & := \Hm_n\left( \tot_{\Pi}\left(\sigma_{\mathrm{I}, \leq 1} \mathrm{CC}(R)\right) \right) & \text{(homologi siklik negatif)}
	\end{align*}
	Keduanya merupakan funktor
	$\Bbbk\dcate{Alg}\to\Bbbk\dcate{Mod}$.
\end{definisi}

% segment-id: o014.li2.chapter3.exercises-hochschild-homology-and-cohomology.p017
Karena $\sigma_{\mathrm{I},\geq0}\mathrm{CC}(R)$ terletak di kuadran pertama,
kompleks totalnya hanya melibatkan jumlah langsung berhingga
$\bigoplus_{\substack{p+q=n\\p,q\geq0}}\mathrm{CC}_{p,q}(R)$. Karena itu,
subskrip pada $\tot$ tidak diperlukan.

% segment-id: o014.li2.chapter3.exercises-hochschild-homology-and-cohomology.p018
Dari definisi langsung diperoleh $\mathrm{HC}_{<0}(R)=0$, sedangkan
$\mathrm{HC}_0(R)$ adalah hasil bagi $R$ oleh
$\Image[d_1:R^{\otimes2}\to R]$ dan
$\Image[\identity-t_0]=0$, yaitu $R/[R,R]$. Bagian latihan akan memberikan
perhitungan untuk lebih banyak kasus khusus.

% segment-id: o014.li2.chapter3.exercises-hochschild-homology-and-cohomology.p019
Kompleks siklik memiliki periodisitas
$\mathrm{CC}(R)=\mathrm{CC}(R)_{\mathrm{I}}[-2]$, yang menghasilkan
isomorfisme $\mathrm{HP}_n(R)\rightiso\mathrm{HP}_{n-2}(R)$. Untuk homologi
siklik, situasi yang bersesuaian ialah epimorfisme ``geser dua kolom ke kiri''
% segment-id: o014.li2.chapter3.exercises-hochschild-homology-and-cohomology.q023
\begin{align*}
	\sigma_{\mathrm{I},\geq0}\mathrm{CC}(R)
	&\to\left(\sigma_{\mathrm{I},\geq0}\mathrm{CC}(R)\right)_{\mathrm{I}}[-2]\\
	\mathrm{CC}(R)_{p,q}
	&\to\begin{cases}
		\mathrm{CC}(R)_{p-2,q}\;\text{(identitas)},&p\geq2,\\
		0,&0\leq p<2.
	\end{cases}
\end{align*}
Kernelnya diberikan oleh kolom ke-$0$ dan ke-$1$ dari $\mathrm{CC}(R)$,
yaitu subbikompleks $\sigma_{\mathrm{I},[0,1]}\mathrm{CC}(R)$. Ini
mendefinisikan operator periodisitas Connes
$S:\mathrm{HC}_n(R)\to\mathrm{HC}_{n-2}(R)$ pada homologi siklik.
\index{operator periodisitas Connes (Connes periodicity operator)}

% segment-id: o014.li2.chapter3.exercises-hochschild-homology-and-cohomology.th001
\begin{teorema}[A.\ Connes]\label{prop:Connes-exact}
	Misalkan $R$ adalah aljabar-$\Bbbk$. Terdapat barisan eksak panjang
	kanonik
% segment-id: o014.li2.chapter3.exercises-hochschild-homology-and-cohomology.q024
	\[ \cdots \xrightarrow{S} \mathrm{HC}_{n-1}(R)
		\xrightarrow{B} \HHm_n(R) \xrightarrow{I} \mathrm{HC}_n(R)
		\xrightarrow{S} \mathrm{HC}_{n-2}(R) \to \cdots, \]
	dengan $S$ operator periodisitas Connes, sedangkan $I$ diinduksi oleh
	penyertaan
	\[
		C_\bullet(R,R)\xrightarrow{\text{kolom ke-$0$}}
		\sigma_{\mathrm{I},\geq0}\mathrm{CC}(R).
	\]
	Adapun $B$ merupakan homomorfisme penghubung yang sesuai.
\end{teorema}

% segment-id: o014.li2.chapter3.exercises-hochschild-homology-and-cohomology.pf002
\begin{bukti}
	Langkah pertama ialah menerapkan
	\eqref{eqn:homology-I-brute-truncation} untuk memperoleh barisan eksak
	pendek bikompleks
% segment-id: o014.li2.chapter3.exercises-hochschild-homology-and-cohomology.g008
	\[\begin{tikzcd}
		0 \arrow[r] & \text{kolom ke-$0$} \arrow[r]
		& \sigma_{\mathrm{I},[0,1]}\mathrm{CC}(R) \arrow[r]
		& \text{kolom ke-$1$} \arrow[r] & 0 .
	\end{tikzcd}\]
	Pengambilan kompleks total hanya melibatkan jumlah langsung berhingga.
	Dengan memeriksanya derajat demi derajat, hasilnya tetap merupakan
	barisan eksak pendek
% segment-id: o014.li2.chapter3.exercises-hochschild-homology-and-cohomology.g009
	\[\begin{tikzcd}
		0 \arrow[r] & \text{kolom ke-$0$} \arrow[r]
		& \tot\left(\sigma_{\mathrm{I},[0,1]}\mathrm{CC}(R)\right)
		\arrow[r] & \text{kolom ke-$1$} \arrow[r] & 0.\\
		& C_\bullet(R,R) \arrow[u,"\sim" sloped]
		& & \mathsf{B}'R \arrow[u,"\sim" sloped] &
	\end{tikzcd}\]
	Kita telah mengetahui bahwa $\mathsf{B}'R$ eksak. Dengan menerapkan
	barisan eksak panjang dari
	Proposisi~\ref{prop:long-exact-sequence-ses}, diperoleh bahwa
	$C_\bullet(R,R)\to
	\tot\left(\sigma_{\mathrm{I},[0,1]}\mathrm{CC}(R)\right)$ merupakan
	kuasi-isomorfisme antarkompleks.

	Selanjutnya, tinjau barisan eksak pendek yang disebutkan ketika
	mendefinisikan $S$:
% segment-id: o014.li2.chapter3.exercises-hochschild-homology-and-cohomology.g010
	\[\begin{tikzcd}
		0 \arrow[r] & \sigma_{\mathrm{I},[0,1]}\mathrm{CC}(R) \arrow[r]
		& \sigma_{\mathrm{I},\geq0}\mathrm{CC}(R) \arrow[r]
		& \left(\sigma_{\mathrm{I},\geq0}\mathrm{CC}(R)\right)_{\mathrm{I}}[-2]
		\arrow[r] & 0 .
	\end{tikzcd}\]
	Setelah diambil kompleks totalnya, barisan ini tetap eksak pendek.
	Berdasarkan langkah sebelumnya, homologi derajat $n$ dari ketiga kompleks
	total itu masing-masing diidentifikasikan dengan
	$\HHm_n(R)$, $\mathrm{HC}_n(R)$, dan $\mathrm{HC}_{n-2}(R)$.\footnote{
	Catatan penerjemah (O014-C048): sumber menyebut ``kohomologi derajat $n$'' pada
	kalimat ini, tetapi ketiga objek adalah kompleks rantai dan rumusnya
	memakai subskrip homologi; ``homologi'' digunakan di sini.}
	Pernyataan teorema pun terbukti.
\end{bukti}

% segment-id: o014.li2.chapter3.exercises-hochschild-homology-and-cohomology.re002
\begin{catatan}
	\index{bikompleks siklik}
	\index{kohomologi siklik periodik (periodic cyclic cohomology)}
	\index{kohomologi siklik (cyclic cohomology)}
	\index[sym1]{HPHC}
	Definisikan bikompleks siklik dalam pengertian kompleks kokrantai,
	$\mathrm{CC}'(R)$, dengan
	$\mathrm{CC}'(R)^{p,q}:=
	\Hom_{\Bbbk}\left(\mathrm{CC}(R)_{p,q},\Bbbk\right)$. Dengan meniru
	Definisi~\ref{def:HC}, definisikan
% segment-id: o014.li2.chapter3.exercises-hochschild-homology-and-cohomology.q025
	\begin{align*}
		\mathrm{HP}^n(R)
		&:=\Hm^n\left(\tot_{\oplus}(\mathrm{CC}'(R))\right)
		&\text{(kohomologi siklik periodik)},\\
		\mathrm{HC}^n(R)
		&:=\Hm^n\left(
		\tot\left(\sigma_{\mathrm{I}}^{\geq0}\mathrm{CC}'(R)\right)\right)
		&\text{(kohomologi siklik)}
	\end{align*}
	dan seterusnya, dengan $\sigma_{\mathrm{I}}^{\geq0}$ tetap menyatakan
	funktor pemenggalan kasar horizontal.

	Tuliskan $R^\vee:=\Hom_{\Bbbk}(R,\Bbbk)$ dan jadikan ini bimodul $(R,R)$
	melalui
	\[ (r\varphi r')(x):=\varphi(r'xr). \]
	Melalui currying, dual $\Bbbk$-linear dari
	kolom Hochschild diidentifikasikan dengan
	$C^\bullet(R,R^\vee)$; identifikasi ini tidak memerlukan hipotesis
	keterhinggaan. Karena itu, barisan eksak panjang pada
	Teorema~\ref{prop:Connes-exact} mempunyai versi kohomologi
% segment-id: o014.li2.chapter3.exercises-hochschild-homology-and-cohomology.q026
	\[ \cdots \xrightarrow{S} \mathrm{HC}^{n+1}(R)
		\xrightarrow{I} \HHm^{n+1}(R,R^\vee) \xrightarrow{B}
		\mathrm{HC}^n(R) \xrightarrow{S} \mathrm{HC}^{n+2}(R) \to \cdots. \]
	\footnote{Catatan penerjemah (O014-C049): sumber mencetak
	$\HHm^{n+1}(R)=\HHm^{n+1}(R,R)$ sebagai suku tengah. Namun dual
	$\Bbbk$-linear dari $C_\bullet(R,R)$ mempunyai koefisien
	$R^\vee$, bukan bimodul reguler $R$, tanpa suatu identifikasi tambahan
	$R^\vee\simeq R$.}
	Operator periodisitas Connes $S$ yang bersesuaian berasal dari
	monomorfisme ``geser dua kolom ke kanan'', yaitu
% segment-id: o014.li2.chapter3.exercises-hochschild-homology-and-cohomology.q027
	\[ \sigma_{\mathrm{I}}^{\geq0}\mathrm{CC}'(R)
		\to\left(\sigma_{\mathrm{I}}^{\geq0}\mathrm{CC}'(R)\right)_{\mathrm{I}}[2]. \]
\end{catatan}

% segment-id: o014.li2.chapter3.exercises-hochschild-homology-and-cohomology.p020
Seiring bertambahnya perangkat yang dikuasai, kita akan berulang kali kembali
membahas homologi dan kohomologi Hochschild serta homologi dan kohomologi
siklik.
